\section{Sự cố và cách xử lý}\label{sec:suco}
Nguyên tắc: \key{bình tĩnh và nói thật}, vừa sửa vừa nói cho thầy biết mình đang làm gì. Quá 30 giây thì bỏ qua và mở ảnh chụp: showcase có ảnh của mọi sơ đồ (mục cuối ``All diagrams'') và mọi màn hình demo (mục F); ảnh gốc ở \file{figures/vp/} và \file{figures/demo/}. Nhiều sự cố của AI lại là dịp chứng minh đường lui (NFR02).

\begin{longtable}{@{}P{3.6cm}P{3.6cm}P{8.6cm}@{}}
\toprule
\textbf{Triệu chứng} & \textbf{Nguyên nhân thường gặp} & \textbf{Cách xử lý}\\
\midrule\endfirsthead
\toprule
\textbf{Triệu chứng} & \textbf{Nguyên nhân thường gặp} & \textbf{Cách xử lý}\\
\midrule\endhead
Showcase không mở được ở \code{localhost:8000} & Terminal 3 chưa chạy, hoặc cổng bận & Chạy lại \code{python3 -m http.server 8001 --bind 127.0.0.1} và mở cổng 8001. Hoặc mở bản trên GitHub Pages (cần Internet).\\
\addlinespace
Không kết nối được \code{127.0.0.1:8686} & Cửa hàng chưa chạy hoặc đã tắt & Trong \file{demo} chạy \code{./run.sh}; kiểm tra bằng \code{lsof -nP -iTCP:8686 -sTCP:LISTEN}.\\
\addlinespace
\emph{Address already in use} & Tiến trình cũ còn chạy & \code{lsof -tiTCP:8686 -sTCP:LISTEN} lấy mã tiến trình, \code{kill <mã>}, rồi \code{./run.sh}.\\
\addlinespace
\code{Gemma 4: offline}, hoặc nhãn ``KEYWORD FALLBACK'' & Mất Internet; key sai hoặc thiếu; dịch vụ Gemma lỗi & Làm theo mục~\ref{sec:gemma}: kiểm tra mạng và key, sửa xong thì khởi động lại cửa hàng. Không sửa được thì trình bày dự phòng có chủ đích (mục~\ref{sec:khonggemma}).\\
\addlinespace
Gemma chậm & Mạng yếu & Mỗi lần gọi chờ tối đa 15 giây rồi tự chuyển sang đường lui. Trong lúc chờ, kể tiếp điều đang xảy ra phía sau.\\
\addlinespace
\emph{Voice search is unavailable} & Whisper chưa chạy & \code{cd demo/voice-server && ./run_voice.sh}, đợi khoảng 10 giây, kiểm tra \code{curl http://127.0.0.1:9000/health}. Tạm thời gõ câu thay vì nói.\\
\addlinespace
\emph{We couldn't hear you}, hoặc nghe sai & Micro tắt; tiếng ồn & Cho phép micro (ổ khóa trên thanh địa chỉ Chrome); nói to, chậm, bằng tiếng Anh; sửa chữ trong ô \emph{We heard}.\\
\addlinespace
Tải ảnh không được & Không phải JPEG, PNG, WebP, hoặc quá 5 MB & Dùng \file{Q02.jpg}, \file{room1.jpg} hoặc ba phòng mẫu có sẵn. Nói luôn: đây là NFR03, máy chủ kiểm tra lại cả khi trình duyệt đã kiểm.\\
\addlinespace
Hết hàng, giỏ lạ, ``latest order'' là đơn cũ & Dữ liệu còn từ lần tập & Đặt lại dữ liệu mẫu (mục~\ref{sec:reset}).\\
\addlinespace
VP Community báo sơ đồ \emph{not supported} & Đó là một trong năm sơ đồ BPMN, ArchiMate, DFD & Bình thường. Mở ảnh chụp trong bài hay showcase (câu trả lời ở mục~\ref{sec:qtrial}); bản dùng thử không mở lại được project.\\
\addlinespace
VP Community không mở được project & Bản quyền, hoặc project đang bị một tiến trình dựng sơ đồ giữ (\file{vp/.vp.lock}) & Đóng mọi cửa sổ VP, xóa \file{vp/.vp.lock} nếu không còn tiến trình nào chạy, mở lại bằng bản Community. Không chạy \file{tools/build_vp.py} trong buổi bảo vệ. Nếu vẫn không được: dùng ảnh trong showcase.\\
\addlinespace
VP hỏi có lưu không & Đã lỡ kéo một hình & Chọn \emph{Don't Save}: file project phải giữ nguyên như lúc nộp.\\
\addlinespace
\code{sqlite_web: command not found} & Chưa cài sqlite-web & Dùng \code{sqlite3 -readonly -header -column demo/data/ecom.db} và các câu ở mục~\ref{sec:cuaso}.\\
\addlinespace
\code{./run.sh} báo thiếu JDK & Chưa cài Java 17 trở lên & \code{brew install openjdk@21}, rồi \code{./build.sh} trong \file{demo}.\\
\addlinespace
Chữ trên máy chiếu nhỏ & Độ phân giải & Phóng to Chrome (Cmd và dấu +) lên 125\%; VS Code: Cmd và dấu +; VP: thanh zoom ở góc dưới.\\
\addlinespace
Mọi thứ hỏng & Máy hoặc mạng & Chỉ cần trình duyệt: showcase (bản trên máy hoặc GitHub Pages) có đủ sơ đồ và ảnh demo. Kể lại năm flow bằng ảnh ở mục F, và nói thẳng rằng demo trực tiếp không chạy được.\\
\bottomrule
\end{longtable}

\begin{note}[warn]{Những điều không nên nói}
\begin{itemize}
\item Không nói Gemma \emph{hơn} luật từ khóa khi tìm bằng chữ: số đo cho thấy ngang. Chỉ nói Gemma hơn rõ ở ảnh.
\item Không nói số đo ở Chương 3 là đo trên DeepInfra: chúng đo ngày 30/10 trên máy chủ của lớp.
\item Không nói VP sinh ra \emph{toàn bộ} cơ sở dữ liệu: bản của VP thiếu khóa ngoại và CHECK, em bổ sung từ cùng mô hình.
\item Không nói mọi yêu cầu đều có kiểm thử tự động; không nói vector ảnh là embedding học được; không nói thanh toán là thật; không nói đặt hàng nằm trong một giao dịch.
\item Không mở \file{demo/config/app.local.properties} khi đang chiếu (có key); không đọc key hay địa chỉ máy chủ lên.
\end{itemize}
\end{note}
